\documentclass[11pt]{article}

\usepackage[a4paper,margin=0.72in]{geometry}
\usepackage{amsmath,amssymb}
\usepackage{booktabs}
\usepackage{tabularx}
\usepackage{array}
\usepackage{enumitem}
\usepackage{xcolor}
\usepackage{microtype}
\usepackage{fancyhdr}
\usepackage{listings}

\pagestyle{fancy}
\fancyhf{}
\lhead{DA3410}
\rhead{Viva Question Set-2}
\cfoot{\thepage}

\setlength{\parindent}{0pt}
\setlength{\parskip}{4pt}
\setlist[itemize]{leftmargin=1.4em,itemsep=1pt,topsep=2pt}
\setlist[enumerate]{leftmargin=1.6em,itemsep=1pt,topsep=2pt}

\lstset{
  language=Python,
  basicstyle=\ttfamily\small,
  columns=fullflexible,
  keepspaces=true,
  showstringspaces=false,
  frame=single,
  framerule=0.3pt,
  aboveskip=4pt,
  belowskip=4pt,
  xleftmargin=0.5em,
  xrightmargin=0.5em
}

% Student copy: answers and marking rubrics are intentionally omitted
% from the source, not merely hidden by a toggle.
\newcounter{question}
\newenvironment{question}{%
  \refstepcounter{question}
  \par\medskip
  \noindent\textbf{Q\thequestion.}\par\smallskip
}{\par\medskip}

\newcommand{}{\textit{Student answers orally; no code writing required.}}

\begin{document}

\begin{center}
    {\LARGE \textbf{DA3410 Viva Questions Set-2}}\\[4pt]
\end{center}

% Bank Q6
\begin{question}
There is one evaluation bug in the code. Identify it. 

\begin{lstlisting}
X = scaler.fit_transform(X)
X_train, X_val, y_train, y_val = train_test_split(X, y)
model.fit(X_train, y_train)
\end{lstlisting}
\end{question}

% Bank Q8
\begin{question}
Give the activation output and local derivative for the input below.
Use the convention that the ReLU derivative at zero is $0$.

\begin{center}
\begin{tabular}{lc}
\toprule
Activation and input & Required output \\
\midrule
ReLU, $x=0$ & $f(x)$ and $f'(x)$ \\
\bottomrule
\end{tabular}
\end{center}
\end{question}

% Bank Q10
\begin{question}
Consider the implementation below.

\begin{verbatim}
h = torch.relu(x @ W1 + b1)
y = h @ W2 + b2
\end{verbatim}

Can the complete mapping from \texttt{x} to \texttt{y} always be replaced by
a single affine layer of the form

\begin{verbatim}
y = x @ W + b
\end{verbatim}

without changing the function? Explain briefly.
\end{question}

% Bank Q14
\begin{question}
Is the checkpointing code correct? If not, identify the problem and state the
fix conceptually.

\begin{verbatim}
best_loss = float("inf")

for epoch in range(epochs):
    ...
    if train_loss < best_loss:
        best_loss = train_loss
        best_params = clone_params(params)
\end{verbatim}
\end{question}

% Bank Q19
\begin{question}
Consider the following manual autodiff forward pass:

\begin{verbatim}
a = multiply(x, w)
b = relu(a)
loss = square(b)
\end{verbatim}

The operations are recorded during the forward pass as

\begin{verbatim}
ops = [multiply_op, relu_op, square_op]
\end{verbatim}

Is the proposed backward implementation correct? If not, state the required
change.

\begin{verbatim}
for op in reversed(ops):
    op.backward()
    op.parents[0].grad = op.grad
\end{verbatim}
\end{question}

% Bank Q23
\begin{question}
Why is the bias gradient reduced with \texttt{sum(dim=0)} rather than
\texttt{sum(dim=1)}? What shape does it have? 

\begin{lstlisting}
grad_z = torch.randn(32, 128)
grad_b = grad_z.sum(dim=0)
\end{lstlisting}
\end{question}

% Bank Q26
\begin{question}
How many batches are produced and what is the size of the last batch?


\begin{lstlisting}
for i in range(0, N, batch_size):
    xb = X[i:i + batch_size]
\end{lstlisting}

\begin{center}
\begin{tabular}{cc}
\toprule
$N$ & \texttt{batch\_size} \\
\midrule
10 & 4 \\
\bottomrule
\end{tabular}
\end{center}
\end{question}

% Bank Q28
\begin{question}
Rank SGD, mini-batch gradient descent, and full-batch gradient descent from
most update noise to least update noise.
\end{question}

% Bank Q30
\begin{question}
A student computes the gradient for the last layer, immediately updates that
layer's weights, and then uses the \emph{updated} weights to propagate gradients
into earlier layers. What is wrong with this procedure?
\end{question}

% Bank Q32
\begin{question}
Suppose \texttt{group\_id[i]} identifies which related group sample
\texttt{i} belongs to. Samples with the same \texttt{group\_id} should not
appear in both training and validation.

Does the code guarantee this? If not, explain the problem.

\begin{verbatim}
train_idx = torch.where(group_id != 5)[0]
val_idx   = torch.where(group_id == 5)[0]
\end{verbatim}
\end{question}


\end{document}
